\chapter{Conclusiones y Trabajo Futuro}
\label{chap:conclusiones}

Este capítulo final sintetiza las contribuciones principales de este Trabajo de Fin de Grado, discute sus implicaciones para el sector \textit{retail} y propone líneas de investigación para continuar evolucionando el sistema.

\section{Conclusiones Generales}
El desarrollo de este sistema de \textit{forecasting} probabilístico ha permitido validar empíricamente que la excelencia en la previsión de demanda no debe medirse únicamente por la reducción del error medio absoluto, sino por la capacidad del modelo para inducir decisiones económicamente óptimas en un entorno de incertidumbre.

Las conclusiones principales se resumen en los siguientes puntos:
\begin{itemize}
    \item \textbf{Desconexión entre error y coste}: Se ha demostrado que modelos con mayor error predictivo puntual (como CatBoost frente a Seasonal Naive) pueden generar ahorros operativos significativos (9.5\%) al gestionar mejor la incertidumbre y la asimetría de costes.
    \item \textbf{Importancia de la Demanda Latente}: La recuperación de la demanda censurada por \textit{stockouts} ha resultado ser el factor más determinante para la rentabilidad, logrando reducciones de coste de hasta el 32\% frente a modelos entrenados con demanda observada.
    \item \textbf{Rigor en la Validación}: El uso de un esquema de validación \textit{walk-forward} estricto ha permitido obtener estimaciones de rendimiento realistas, evitando el sobreajuste y el \textit{leakage} temporal que invalidarían los resultados académicos.
    \item \textbf{Desafíos en la Calibración}: A pesar del éxito económico, la sub-cobertura detectada en los intervalos de confianza revela que la cuantificación de la incertidumbre aleatoria sigue siendo un reto complejo que requiere técnicas de calibración proactivas.
\end{itemize}

\section{Contribución al Desarrollo Sostenible}
Alineado con el \textbf{ODS 12 (Producción y Consumo Responsables)}, este trabajo aporta una herramienta técnica capaz de mitigar el desperdicio alimentario. Al optimizar las órdenes de compra mediante el modelo del \textit{Newsvendor}, el sistema reduce sistemáticamente el sobrestock de productos perecederos, disminuyendo la huella ambiental de la cadena de suministro sin comprometer la seguridad alimentaria ni el nivel de servicio al cliente.

\section{Líneas de Trabajo Futuro}
La arquitectura modular desarrollada deja el camino preparado para futuras expansiones:
\begin{itemize}
    \item \textbf{Modelos de Deep Learning}: Explorar el uso de redes neuronales temporales (ej. LSTM o Transformers) dentro de la arquitectura global para capturar dependencias de largo plazo más complejas.
    \item \textbf{Calibración Avanzada}: Implementar métodos de \textit{Conformal Prediction} dinámicos que se ajusten en tiempo real según la variabilidad local de cada tienda o categoría de producto.
    \item \textbf{Compensación de Stockout Multietapa}: Desarrollar algoritmos de imputación más sofisticados que consideren el efecto de sustitución entre productos cuando ocurre una rotura de stock.
    \item \textbf{Integración con Logística}: Evolucionar de un modelo de una sola etapa (\textit{Newsvendor}) a una política de inventario multietapa que considere el \textit{lead time} variable y los costes de transporte.
\end{itemize}

Como cierre, este TFG demuestra que la integración de la estadística avanzada con la ingeniería de software permite transformar datos brutos en valor estratégico, proporcionando una base sólida para la digitalización inteligente del sector \textit{retail}.
